\section{Conclusion}

This paper documents a clear firm-size wage premium in Colombia using harmonized GEIH microdata for 2008--2019 and 2021--2025. The descriptive evidence shows that real hourly wages rise sharply with firm size, while employment remains concentrated in solo work and micro-firms. Solo workers and workers in firms with up to ten workers account for about two thirds of employment, and the middle of the firm-size distribution remains thin.

The regression results show that a substantial conditional premium remains after controlling for gender, age, age squared, education, formality, and sector-year fixed effects. In the full specification, workers in firms with 101 or more employees earn about 43 percent more than otherwise comparable solo workers. The premium is positive throughout the sample period. The point estimates suggest an increase between 2008 and 2025, but the estimated change is not statistically significant.

The formal-informal results sharpen the main finding. Formal workers earn more than informal workers within every firm-size category. At the same time, the firm-size premium is much steeper among informal workers than among formal workers. Within formal employment, the premium relative to formal solo workers is small, although workers in 101+ firms earn more than workers in intermediate-size formal firms. Taken together, these results show that the Colombian firm-size wage premium is
not only a formal-informal contrast. It is a persistent wage gap associated with firm scale, especially within informal employment, and it remains after conditioning on the main observable worker and job characteristics available in
the GEIH.

%The approximate elasticity estimates reinforce this conclusion. The full-sample conditional elasticity is 0.071, above the observable-controls benchmark in the meta-analysis of \citet{DiegmannMullerSchoefer2026}, but the split by formality shows that this relatively high elasticity is mainly an informal-sector phenomenon: the estimate is 0.026 among formal workers and 0.130 among informal workers.
